% ============================================================
% PREÁMBULO — Configuración global del documento
% ============================================================
%
% CONTENIDO:
%   1. Presentación     — fondo de carátula (background)
%   2. Codificación     — inputenc (UTF-8), fontenc (T1), lmodern, courier
%   3. Matemáticas      — amsmath, amssymb, amsfonts, latexsym, textcomp, cancel
%   4. Tablas           — array, tabularx, multirow, longtable, booktabs
%   5. Figuras          — graphicx, caption, subfigure, float, rotating
%   6. Colores          — xcolor + paletas en Colores.tex
%   7. Listados         — listings + estilos matlab_dark, asmstyle, miasm
%   8. Cajas            — tcolorbox
%   9. TikZ             — tikz + librerías + estilos en TikzEstilos.tex
%  10. Gráficas         — pgfplots, circuitikz
%  11. Página           — geometry, fancyhdr, titlesec
%  12. Listas           — enumerate, paralist, easylist, enumitem
%  13. Extras           — pdfpages, verbatim, parindent, hyphenation
%  14. Comandos         — nombres en castellano, atajos de bloques
%
% ARCHIVOS ASOCIADOS (cargados al final de cada sección):
%   Colores.tex         — todas las paletas de color
%   TikzEstilos.tex     — \tikzset para bloques y diagramas de flujo
%
% ============================================================

% ============================================================
% MULTI-ARCHIVO
% ============================================================

% --- Permite compilar cada capítulo de forma independiente ---
\usepackage{subfiles}

% ============================================================
% PRESENTACIÓN
% ============================================================

% --- Imagen de fondo en la carátula ---
\usepackage{background}

% ============================================================
% CODIFICACIÓN Y FUENTES
% ============================================================

% --- Codificación UTF-8 y separación correcta de palabras con acentos ---
\usepackage[utf8]{inputenc}
% --- Fuente de salida con soporte completo de caracteres T1 ---
\usepackage[T1]{fontenc}
% --- Fuente vectorial moderna (reemplaza Computer Modern bitmap) ---
\usepackage{lmodern}
% --- Fuente monoespaciada Courier (para código) ---
\usepackage{courier}

% ============================================================
% MATEMÁTICAS
% ============================================================

% --- Entornos y símbolos matemáticos avanzados (AMS) ---
\usepackage{amsmath,amssymb,amsfonts}
% --- Símbolos extra: latexsym, textcomp, cancelación de términos ---
\usepackage{latexsym,textcomp,cancel}

% ============================================================
% TABLAS
% ============================================================

% --- Columnas con formato avanzado (p, m, b, >{}, etc.) ---
\usepackage{array}
% --- Tablas con ancho ajustable al texto ---
\usepackage{tabularx}
% --- Celdas que abarcan varias filas ---
\usepackage{multirow}
% --- Tablas que se extienden a lo largo de varias páginas ---
\usepackage{longtable}
% --- Líneas horizontales profesionales (\toprule, \midrule, \bottomrule) ---
\usepackage{booktabs}
% --- Colores en celdas y filas de tablas ---
\usepackage{colortbl}

% ============================================================
% FIGURAS Y GRÁFICOS
% ============================================================

% --- Inserción de imágenes externas ---
\usepackage{graphicx}
\DeclareGraphicsExtensions{.pdf,.png,.jpg}
\graphicspath{{Graficos/}}

% --- Control avanzado de leyendas (caption) ---
\usepackage{caption}
% --- Subfiguras dentro de una figura ---
\usepackage{subfigure}
% --- Posicionamiento forzado de flotantes con [H] ---
\usepackage{float}
% --- Rotación de figuras y tablas ---
\usepackage{rotating}

% ============================================================
% COLORES
% ============================================================

% --- Motor de color con soporte para tablas y nombres extendidos ---
\usepackage[table,dvipsnames]{xcolor}

% --- Paletas de color (matlab, flujo, acento, tablas, informes) ---
% ============================
% COLORES — Colores.tex
% ============================
%
% CONTENIDO:
%   1. matlab_dark  — colores para el estilo de listados MATLAB
%   2. Flujo        — paleta completa para diagramas de flujo TikZ
%   3. Acento       — paleta c1–c9 (azul oscuro → rojo)
%   4. Tablas       — encabezado y filas alternas
%   5. Informes     — cajas, texto secundario, indicadores
%
% NOTA: Requiere \usepackage[table,dvipsnames]{xcolor} — cargado en Preambulo.tex
%
% ============================

% ── Colores para el estilo matlab_dark ──────────────────────
\definecolor{bg_matlab}{RGB}{235, 242, 252}   % azul muy claro (fondo)
\definecolor{matlab_kw}{RGB}{0,   60, 160}    % azul navy (keywords)
\definecolor{matlab_str}{RGB}{0,  120,  60}   % verde (strings)
\definecolor{matlab_com}{RGB}{110, 110, 110}  % gris (comentarios)
\definecolor{matlab_num}{RGB}{110, 110, 110}  % gris (números de línea)

% ── Colores para diagramas de flujo ─────────────────────────
\definecolor{vacioQ}{RGB}{255, 252, 242}
\definecolor{vacioBQ}{RGB}{180, 180, 200}
\definecolor{blancoW}{RGB}{255, 255, 255}
\definecolor{blancoBW}{RGB}{200, 200, 200}
\definecolor{poloE}{RGB}{220, 220, 220}
\definecolor{poloBE}{RGB}{100, 100, 100}
\definecolor{negroR}{RGB}{30,  30,  30}
\definecolor{negroBR}{RGB}{0,   0,   0}
\definecolor{cafeT}{RGB}{210, 190, 170}
\definecolor{cafeBT}{RGB}{130, 100,  80}
\definecolor{azulA}{RGB}{150, 230, 240}
\definecolor{azulBA}{RGB}{0,  120, 150}
\definecolor{azulS}{RGB}{150, 200, 255}
\definecolor{azulBS}{RGB}{0,  100, 200}
\definecolor{lilaD}{RGB}{190, 170, 250}
\definecolor{lilaBD}{RGB}{100,  70, 200}
\definecolor{rojoG}{RGB}{255, 160, 180}
\definecolor{rojoBG}{RGB}{200,   0,   0}
\definecolor{verdeZ}{RGB}{170, 240, 190}
\definecolor{verdeBZ}{RGB}{0,  150,   0}
\definecolor{verdeX}{RGB}{150, 255, 220}
\definecolor{verdeBX}{RGB}{0,  150, 120}
\definecolor{amarilloC}{RGB}{255, 230, 150}
\definecolor{amarilloBC}{RGB}{230, 150,   0}
\definecolor{naranjaV}{RGB}{255, 200, 150}
\definecolor{naranjaBV}{RGB}{220, 100,   0}
\definecolor{rojoB}{RGB}{255, 150, 150}
\definecolor{rojoBB}{RGB}{200,   0,   0}
\definecolor{rosaF}{RGB}{255, 200, 220}
\definecolor{rosaBF}{RGB}{200,  80, 130}
\definecolor{darkgreen}{RGB}{0,  130,   0}

% ── Paleta de acento (c1–c9) ────────────────────────────────
\definecolor{c1}{HTML}{011627}
\definecolor{c2}{HTML}{004E67}
\definecolor{c3}{HTML}{008891}
\definecolor{c4}{HTML}{9FD9C3}
\definecolor{c5}{HTML}{E6D7A7}
\definecolor{c6}{HTML}{F6A600}
\definecolor{c7}{HTML}{D2691E}
\definecolor{c8}{HTML}{A33300}
\definecolor{c9}{HTML}{7C1E18}

% ── Colores para tablas con rowcolor ────────────────────────
\definecolor{azulprinc}{RGB}{46, 116, 181}   % encabezado de tabla
\definecolor{grisclaro}{RGB}{242, 242, 242}  % filas alternas

% ── Colores generales de informes ────────────────────────────
\definecolor{azulclaro}{RGB}{210, 225, 245}  % fondo de cajas y resaltados
\definecolor{textgris}{RGB}{90, 90, 90}      % texto secundario (encabezados)
\definecolor{verdeclaro}{RGB}{220, 240, 220} % fondo indicador positivo
\definecolor{rojoclaro}{RGB}{250, 225, 225}  % fondo indicador negativo


% ============================================================
% LISTADOS DE CÓDIGO
% ============================================================

% --- Inserción y resaltado de código fuente ---
\usepackage{listings}

% ── Estilo global por defecto ────────────────────────────────
\lstset{
  basicstyle=\scriptsize\ttfamily,
  numbers=left,
  numberstyle=\tiny\color{matlab_com},
  stepnumber=1,
  numbersep=5pt,
  breaklines=true,
  captionpos=b,
  showspaces=false,
  showstringspaces=false,
  showtabs=false,
  tabsize=4
}

% ── Estilo MATLAB (fondo azul pálido, bajo consumo de tinta) ─
\lstdefinestyle{matlab_dark}{
  language=Matlab,
  captionpos=b,
  backgroundcolor=\color{bg_matlab},
  basicstyle=\ttfamily\small\color{black},
  keywordstyle=\color{matlab_kw}\bfseries,
  stringstyle=\color{matlab_str},
  commentstyle=\color{matlab_com}\itshape,
  numberstyle=\tiny\color{matlab_num},
  numbers=left,
  numbersep=10pt,
  frame=single,
  framerule=0pt,
  rulecolor=\color{bg_matlab},
  breaklines=true,
  tabsize=4,
  showstringspaces=false,
  xleftmargin=16pt,
  framexleftmargin=16pt
}

% ── Estilo ensamblador x86 ───────────────────────────────────
\lstdefinestyle{asmstyle}{
  language=[x86masm]Assembler,
  backgroundcolor=\color{white},
  basicstyle=\ttfamily\footnotesize,
  keywordstyle=\color{magenta},
  commentstyle=\color{verdeBZ},
  stringstyle=\color{lilaBD},
  numberstyle=\tiny\color{matlab_com},
  numbers=left,
  numbersep=5pt,
  frame=single,
  breaklines=true,
  captionpos=b,
  keepspaces=true,
  showspaces=false,
  showstringspaces=false,
  showtabs=false,
  tabsize=4
}

% ── Estilo ensamblador genérico (sin keywords fijos) ─────────
\lstdefinestyle{miasm}{
  basicstyle=\ttfamily\footnotesize,
  keywordstyle=\color{blue},
  commentstyle=\color{gray},
  columns=fullflexible,
  keepspaces=true,
  frame=single,
  numbers=left,
  numberstyle=\tiny\color{gray},
  xleftmargin=1em,
  tabsize=4
}

% ============================================================
% CAJAS Y ENTORNOS ESPECIALES
% ============================================================

% --- Cajas con color, bordes y fondos personalizables ---
\usepackage{tcolorbox}

% ============================================================
% TIKZ — GRÁFICOS VECTORIALES
% ============================================================

% --- Motor principal de diagramas vectoriales ---
\usepackage{tikz}
\usetikzlibrary{
  arrows.meta,
  shapes.geometric,
  shapes.misc,
  positioning,
  calc,
  decorations.markings,
  shapes,
  arrows,
  fit
}

% --- Estilos de bloques y diagramas de flujo ---
% ============================
% ESTILOS TIKZ — TikzEstilos.tex
% ============================
%
% CONTENIDO:
%   1. Bloques      — block, gain, port, sum, joint, input, output, box,
%                     lifeline, copyarrow, circ
%   2. Flujo        — terminator, process, decision, io, coment, document,
%                     subprocess, data, database, connector, cloud,
%                     startstop, loop, junction, line*, arrow
%
% NOTA: Requiere \usepackage{tikz} y \usetikzlibrary{...} — cargados en Preambulo.tex
%       Los colores de flujo vienen de Colores.tex
%
% ============================

% ── Estilos genéricos de diagramas de bloques ────────────────
\tikzset{
  block/.style     = {draw, thick, rectangle, minimum height=10mm, minimum width=3em},
  block1/.style    = {draw=none, rectangle, minimum height=7mm, minimum width=3em},
  gain/.style      = {draw, thick, isosceles triangle, isosceles triangle apex angle=45},
  port/.style      = {inner sep=0pt, font=\tiny},
  sum/.style n args={4}{draw, circle, node distance=2cm, minimum size=7mm, alias=sum,
    append after command={
      node at (sum.north) [port, below=1pt] {$#1$}
      node at (sum.west)  [port, right=1pt] {$#2$}
      node at (sum.south) [port, above=1pt] {$#3$}
      node at (sum.east)  [port, left=1pt]  {$#4$}
    }},
  joint/.style     = {circle, draw, fill, inner sep=0pt, minimum size=2pt},
  jointi/.style    = {circle, draw, fill, inner sep=0pt, minimum size=0pt},
  input/.style     = {coordinate},
  output/.style    = {coordinate},
  box/.style       = {rectangle, draw, fill=blue!15, minimum width=3cm,
                      minimum height=1cm, text centered, font=\bfseries},
  lifeline/.style  = {dashed, thick},
  copyarrow/.style = {-{Stealth}, thick},
  circ/.style      = {circle, fill=black, inner sep=2pt}
}

% ── Estilos para diagramas de flujo ─────────────────────────
\tikzset{
  terminator/.style = {ellipse, draw=verdeBZ, fill=verdeZ,
    text centered, minimum height=2em, minimum width=6em},
  process/.style    = {rectangle, draw=azulBS, fill=azulS,
    text width=10em, text centered, rounded corners, minimum height=2em},
  decision/.style   = {diamond, draw=rojoBG, fill=rojoG,
    text width=6em, text badly centered, inner sep=0pt},
  io/.style         = {trapezium, trapezium left angle=70, trapezium right angle=110,
    draw=verdeBX, fill=verdeX, text centered,
    minimum height=2em, minimum width=8em},
  coment/.style     = {rectangle, draw=vacioBQ, fill=vacioQ, dashed,
    text width=15em, align=justify, minimum height=4em, inner sep=6pt},
  document/.style   = {rectangle, draw=azulBA, fill=azulA,
    text centered, minimum height=2em, minimum width=8em,
    append after command={
      (\tikzlastnode.south west) .. controls +(0,-0.5) and +(0,-0.5) ..
      (\tikzlastnode.south east) -- cycle}},
  subprocess/.style = {rectangle, draw=rosaBF, fill=rosaF,
    double, double distance=1pt, text centered,
    minimum height=2em, minimum width=8em},
  data/.style       = {cylinder, draw=amarilloBC, fill=amarilloC,
    shape border rotate=90, aspect=0.25,
    text centered, minimum height=3em, minimum width=4em},
  database/.style   = {cylinder, draw=naranjaBV, fill=naranjaV,
    shape border rotate=90, aspect=0.15,
    text centered, minimum height=3em, minimum width=6em},
  connector/.style  = {circle, draw=negroBR, fill=blancoW,
    minimum size=1.5em, text centered},
  cloud/.style      = {ellipse, draw=cafeBT, fill=cafeT,
    text centered, minimum height=3em, minimum width=6em},
  startstop/.style  = {ellipse, draw=darkgreen, fill=verdeZ,
    minimum width=20mm, minimum height=6mm, align=center, font=\footnotesize},
  loop/.style       = {rectangle, draw=black, fill=yellow!40,
    minimum width=50mm, minimum height=8mm, align=center, font=\footnotesize},
  junction/.style   = {circle, draw=black, fill=black,
    minimum size=4pt, inner sep=0pt},
  line/.style       = {draw=black, -latex', very thick},
  lineRed/.style    = {draw=rojoBG, -latex', very thick},
  lineBlue/.style   = {draw=azulBA, -latex', very thick},
  lineGreen/.style  = {draw=verdeBX, -latex', very thick},
  lineComent/.style = {dashed, draw=vacioBQ},
  arrow/.style      = {-Stealth, thick}
}


% ============================================================
% PGFPLOTS Y CIRCUITIKZ
% ============================================================

% --- Gráficas de funciones y datos con TikZ ---
\usepackage{pgfplots}
\pgfplotsset{compat=1.18}

% --- Diagramas de circuitos eléctricos ---
\usepackage[siunitx]{circuitikz}

% ============================================================
% MÁRGENES Y FORMATO DE PÁGINA
% ============================================================

% --- Márgenes del documento ---
\usepackage[top=2.5cm, bottom=2.5cm, left=3cm, right=2cm]{geometry}
% --- Encabezados y pies de página personalizados ---
\usepackage{fancyhdr}
% --- Formato de títulos de sección (usado en EstiloSecciones.tex) ---
\usepackage{titlesec}

% ============================================================
% LISTAS
% ============================================================

% --- Listas numeradas con estilos personalizables ---
\usepackage{enumerate}
% --- Listas compactas (inparaenum, compactitem, etc.) ---
\usepackage{paralist}
% --- Listas simplificadas con @ como marcador ---
\usepackage{easylist}
% --- Control fino de márgenes en listas (itemsep, leftmargin, etc.) ---
\usepackage{enumitem}

% ============================================================
% EXTRAS
% ============================================================

% --- Inserción de páginas PDF completas ---
\usepackage{pdfpages}
% --- Bloque de texto literal sin interpretación ---
\usepackage{verbatim}

% Sin sangría en párrafos
\parindent=0pt
% Palabras que LaTeX no debe partir con guión
\hyphenation{Linux manera}

% ============================================================
% COMANDOS PERSONALIZADOS
% ============================================================

% Nombres en castellano
\renewcommand{\tablename}{Tabla}
\renewcommand{\figurename}{Figura}
\renewcommand{\contentsname}{Contenido}
\renewcommand{\lstlistingname}{Listado}

% Atajos para diagramas de bloques
\newcommand{\suma}{\Large$+$}
\newcommand{\inte}{$\displaystyle \int$}
\newcommand{\derv}{\huge$\frac{d}{dt}$}
